\chapter{Phụ lục cấu hình thực nghiệm}
\label{app:cau-hinh}

Phụ lục này ghi lại các cấu hình chính dùng trong báo cáo. Mọi tham số thực nghiệm nằm trong
\repoFile{configs/}; không có tham số nào được hard-code trong code. Một cấu hình được ghép theo
thứ tự \texttt{defaults} → nội dung file → override dòng lệnh. Dict ghép đệ quy; list bị thay
nguyên khối. Đường dẫn tương đối tính từ gốc repo.

\section{Cấu hình gốc}
\begin{lstlisting}[numbers=none]
project: caro-11x11
seed: 2026
device: auto
engine: {size: 11, win_length: 5, layout: stride11}
experiment: {kind: misc, name: null, description: "", rq: []}
run: {root: runs, tag: null, index: runs/index.jsonl}
log: {level: INFO, console: true, file: run.log, games: games.jsonl,
      metrics: metrics.jsonl, per_move_stats: true, flush_every: 1}
checkpoint: {dir: checkpoints, keep_last: 5}
\end{lstlisting}
\texttt{seed} là seed gốc; seed từng ván được dẫn xuất bằng
\texttt{derive\_seed(seed, game\_id, color)}. \texttt{device} nhận \texttt{auto}, \texttt{cpu}
hoặc \texttt{cuda}. \texttt{per\_move\_stats} bật ghi thời gian/nút/mô phỏng từng nước; tắt khi
log quá lớn.

\section{Ngân sách}
\begin{lstlisting}[numbers=none]
# mode_a_time.yaml
budget: {mode: time, time_ms: 1000, max_depth: null, max_sims: null}
# mode_b_fixed.yaml
budget: {mode: fixed, time_ms: null, max_depth: 4, max_sims: 800}
\end{lstlisting}
Chế độ A dùng thời gian; chế độ B dùng độ sâu hoặc số mô phỏng. Trong RR8, ngân sách được ghi
đè theo từng tác tử: L3 độ sâu 4, L4 400 mô phỏng, L7/L8 200 mô phỏng. Lý do: một ngân sách
chung 800 mô phỏng cho PUCT sẽ rất chậm, trong khi 200 mô phỏng đã đủ để thấy khác biệt.

\section{Tác tử}
\begin{center}\small
\begin{tabularx}{\textwidth}{@{}l l X@{}}
\toprule
\textbf{Cấp} & \textbf{Kind} & \textbf{Tham số chính} \\
\midrule
L1 & \texttt{random} & không \\
L2 & \texttt{rule\_based} & priorities P1--P7, radius 2, eval\_weights default \\
L3 & \texttt{alphabeta} & pruning, move\_ordering, tt, tt\_size\_mb 64, iterative\_deepening, forced\_moves, radius 2, backend auto \\
L4 & \texttt{mcts} & c\_uct 1.4, rollout random/light, rollout\_radius 1, candidate\_radius 2, root\_tactics \\
L5 & \texttt{policy} & checkpoint model:l5, temperature 0 \\
L6 & \texttt{policy} & checkpoint model:l6, temperature 0 \\
L7 & \texttt{puct} & checkpoint model:l5, c\_puct 1.5, leaf\_eval value, batch\_leaves 8 \\
L8 & \texttt{puct} & checkpoint model:l8/local, c\_puct 1.5, leaf\_eval value, batch\_leaves 8 \\
\bottomrule
\end{tabularx}
\end{center}

\paragraph{Alpha-Beta.} Các cờ \texttt{pruning}, \texttt{move\_ordering}, \texttt{tt},
\texttt{forced\_moves} tồn tại để làm ablation A1/A2. Mặc định bật hết. \texttt{backend}
\texttt{auto} dùng Cython nếu đã build, không thì Python. \texttt{eval\_weights} nhận
\texttt{default}, \texttt{no\_defense} hoặc dict trọng số.

\paragraph{MCTS.} \texttt{rollout} nhận \texttt{random} hoặc \texttt{light}. \texttt{light}
ưu tiên thắng ngay, rồi chặn thua ngay, còn lại ngẫu nhiên trong bán kính 1. \texttt{root\_tactics}
kiểm tra thắng/chặn/đòn kép trước khi mô phỏng. \texttt{reuse\_tree} tái sử dụng cây giữa các
nước; mặc định tắt để đo số mô phỏng công bằng.

\paragraph{PUCT.} \texttt{leaf\_eval} nhận \texttt{value} (chỉ value của mạng) hoặc \texttt{mix}
(trộn value với rollout, kiểu AlphaGo). \texttt{dirichlet} chỉ dùng khi tự chơi; thi đấu đặt
\texttt{null}. \texttt{batch\_leaves} gom lá để suy luận theo lô. \texttt{temperature} 0 khi
thi đấu, $>0$ khi tự chơi.

\section{Huấn luyện}
\paragraph{Sinh dữ liệu.}
\begin{lstlisting}[numbers=none]
datagen:
  expert: {from: agents/l3_alphabeta, params: {backend: auto}}
  expert_depth: 3
  games: 20000
  top_k: 3
  prune_root: false
  temperature: 0.5
  split: [0.8, 0.1, 0.1]
  dedup: d4_canonical
  out: runs/data/expert_v1
\end{lstlisting}
Bản v2 đổi \texttt{temperature: 8.0}, \texttt{top\_k: 5}, \texttt{out: runs/data/expert\_v2}.
Đây là thay đổi tạo ra bước nhảy top-1 từ 44,0\% lên 65,3\%.

\paragraph{SL.}
\begin{lstlisting}[numbers=none]
train: {optimizer: adam, lr: 1.0e-3, weight_decay: 1.0e-4,
        batch_size: 256, epochs: 30, value_coef: 1.0}
data: {dir: runs/data/expert_v1, augment_x8: true}
\end{lstlisting}
Bản v3 dùng \texttt{epochs: 60}, \texttt{lr\_schedule: cosine}, \texttt{data.dir: expert\_v2}.
Nhãn value được đổi sang góc nhìn người đi trong \texttt{train\_sl.py}.

\paragraph{REINFORCE.}
\begin{lstlisting}[numbers=none]
train: {init: model:l5, episodes: 50000, batch_games: 128, lr: 2.0e-4,
        gamma: 1.0, baseline: value_head, entropy_coef: 0.01,
        entropy_decay: 0.999, temperature: 1.0,
        opponent_pool: {latest: 0.5, past_checkpoints: 0.5},
        eval_every: 1000, eval_vs: [agents/l2_rule, agents/l3_alphabeta],
        eval_games: 100}
\end{lstlisting}
Mỗi vòng chơi 128 ván song song, cập nhật một bước, giảm entropy theo hệ số 0,999. Đánh giá mỗi
1\,000 episode, 100 ván trước L2 và L3.

\paragraph{AlphaZero.}
\begin{lstlisting}[numbers=none]
selfplay: {backend: auto, parallel_games: auto, games_per_iter: 200,
           sims_per_move: 400, c_puct: 1.5,
           dirichlet: {alpha: 0.3, epsilon: 0.25}, temperature_moves: 8}
buffer: {window_games: 5000, augment_x8: true}
train: {init: model:l5, optimizer: adam, lr: 1.0e-3, weight_decay: 1.0e-4,
        batch_size: 256, steps_per_iter: 1000, iterations: 50}
gating: {games: 100, swap_colors: true, threshold: 0.55, sims_per_move: 200}
evaluation: {games: 20, sims_per_move: 200}
\end{lstlisting}
Mỗi vòng 200 ván tự đấu ở 400 mô phỏng, 1\,000 bước cập nhật, 100 ván gating ở 200 mô phỏng.
Đánh giá cuối 20 ván trước L7.

\section{Giải đấu}
\begin{lstlisting}[numbers=none]
arena:
  format: round_robin
  games_per_pair: 200
  swap_colors: true
  max_plies: 121
  workers: 4
  openings: {count: 100, plies: [2, 4], region: 5, seed: 7,
             file: runs/openings/openings_v1.json}
  stats: {ci: wilson, confidence: 0.95, bootstrap: 1000, elo: true}
\end{lstlisting}
RR8 dùng \texttt{num\_openings: 50}, \texttt{workers: 4} (hoặc 1 khi chạy nền). Mỗi khai cuộc
2--4 nước trong vùng 5×5 giữa bàn, sinh một lần và dùng chung.

\section{Model}
\begin{lstlisting}[numbers=none]
model: {arch: resnet_pv, in_planes: 4, channels: 64, blocks: 6,
        policy_head: {conv_out: 2, fc_out: 121},
        value_head: {conv_out: 1, fc_hidden: 64}}
\end{lstlisting}
Code thực tế dùng \texttt{value\_fc1} 32 nơ-ron, không phải 64 như file YAML ghi. Khi lệch,
code là chuẩn. Mạng có khoảng 0,48 triệu tham số.

\section{Con trỏ model}
\begin{lstlisting}[numbers=none]
repo: solozy/CaroGame
models:
  l5: {file: l5/best.pt, revision: a76d294b...}
  l6: {file: l6/best.pt, revision: 3ecb6fda...}
  l8: {file: l8/best.pt, revision: null}
\end{lstlisting}
\texttt{model:l5} tải đúng revision đã ghim. L8 chưa push tại thời điểm viết, nên báo cáo dùng
checkpoint cục bộ. Muốn tái lập L8 cần file \texttt{best.pt} của run gốc.

\section{Ghi đè khi chạy}
Ví dụ:
\begin{lstlisting}[numbers=none]
uv run caro-config experiments/rr8_mode_b arena.workers=1
uv run python -m caro.arena.tournament experiments/rr8_mode_b arena.workers=1
uv run python -m caro.training.train_pg training/pg train.init=random run.tag=from_random
\end{lstlisting}
Override dùng dạng \texttt{a.b=c}; giá trị đọc theo YAML. Trong file YAML, số mũ phải có dấu
chấm (\texttt{3.0e-4}), không viết \texttt{3e-4} vì PyYAML đọc thành chuỗi.